\documentclass[margin,line]{res}

\usepackage{hyperref}
\hypersetup{colorlinks=true,urlcolor=blue}

\oddsidemargin -.5in
\evensidemargin -.5in
\textwidth=6.0in
\itemsep=0in
\parsep=0in

\newenvironment{list1}{
  \begin{list}{\ding{113}}{%
      \setlength{\itemsep}{0in}
      \setlength{\parsep}{0in} \setlength{\parskip}{0in}
      \setlength{\topsep}{0in} \setlength{\partopsep}{0in} 
      \setlength{\leftmargin}{0.17in}}}{\end{list}}
\newenvironment{list2}{
  \begin{list}{$\bullet$}{%
      \setlength{\itemsep}{0in}
      \setlength{\parsep}{0in} \setlength{\parskip}{0in}
      \setlength{\topsep}{0in} \setlength{\partopsep}{0in} 
      \setlength{\leftmargin}{0.2in}}}{\end{list}}


\begin{document}
\vspace*{-.5in}
\name{Jason F. Rowe \vspace*{.1in}}

\begin{resume}
\section{\sc Contact Information}
\vspace{.05in}
\begin{tabular}{@{}p{3in}p{4in}}
Bishop's Univeristy                    &  \\         
Dept of Physics and Astronomy    &   {\it Voice:}  (819) 822-9600 x2049  \\         
2600 College Street                    & {\it E-mail:}  Jason.Rowe@ubishops.ca\\
Sherbrooke, QC, Canada, J1M 1Z7        & {\it www:} physics.ubishops.ca/exoplanets \\
\end{tabular}

\vspace*{-.1in}

\section{\sc Employment}
{\bf Full Professor} Bishop's University, July 2021 to present \\
{\bf Associate Professor} Bishop's University, July 2019 to July 2021 \\
{\bf Canada Research Chair in Exoplanet Astrophysics}, July 2017 to present \\
{\bf Assistant Professor} Bishop's University, July 2017 to July 2019\\
%\vspace*{-.2in}
{\bf Principle Investigator}, SETI Institute, November 2011 to present\\
%\vspace*{-.2in}
{\bf JWST NIRISS/FGS Research Scientist}, Universit\'e de Montr\'eal, Sept. 2015 to July 2017\\
%\vspace*{-.2in}
{\bf Kepler Science Office}, NASA-Ames Research Center, November 2010 to Sept. 2015\\
%\vspace*{-.2in}
{\bf NASA Postdoctoral Program}, NASA-Ames Research Center, November 2007 to Sept. 2010%\\

\vspace*{-.1in}

\section{\sc Education}
{\bf University of British Columbia}, Vancouver, British Columbia, Canada\\
\vspace*{-.15in}
\begin{list1}
\item[] Ph.D., Physics and Astronomy, April 2008
\item[] M.Sc., Astronomy,  October 2003
\end{list1}
\vspace*{-.15in}
{\bf University of Toronto}, Toronto, Ontario, Canada\\
\vspace*{-.15in}
\begin{list1}
\item[] B.Sc., Physics and Astronomy,  May, 2000
\end{list1}

\vspace*{-.1in}

\section{\sc Research Interests}
Wide interests in astronomy and astrophysics.  Focused research and studies on the properties, statistics and occurrence of extrasolar planets and the structure and evolution of stars.  Interests and continuing research regarding star clusters, compact objects, microlensing, seismology, habitability and SETI.  Wealth of experience with various space based astronomical instruments and platforms and desire to help develop future missions.  Supporter of \href{https://github.com/jasonfrowe}{Open Source} research software.  

\vspace*{-.1in}

\section{\sc Research and Grants}
{\bf CSA Class G\&C Program to Support Research, Awareness and Learning in Space Science and Technology} 2023-2026 : Pandora Mission\\
{\bf Compute Canada} 2023 - 2027 : Exoplanet Occurrence Rates \\
{\bf Canada Research Chair} 2023 - 2028: Exoplanet Astrophysics (\$500K CAD)\\
{\bf Canadian Space Agency Space Technology Development Program} 2021 - 2022 : Photometric Observations of Transiting Extrasolar Planets. (\$1.1M CAD) \\
{\bf Canadian Space Agency Science Maturation Study} 2018 - 2019 : Photometric Observations of Transiting Extrasolar Planets. (\$150K CAD) \\
{\bf FQRNT New University Researchers Startup Program}  2018 - 2021 : Photometric Observations of Solar System Giants. (\$40K CAD) \\
{\bf NSERC Discovery Grant}  2017 - 2023: Determining What Makes a Planet Earth-like. (\$125K CAD)\\
{\bf CFI Grant} 2017: Determining What Makes a Planet Earth-like (\$400K CAD)\\
%{\bf Canada Research Chair} 2017 - 2022: Exoplanet Astrophysics (\$500K CAD)\\
%{\bf Compute Canada} 2016 - 2022 : Microlensing and Planets in the Kepler FOV 
%{\bf FQRNT Regroupement Strategique} 2017 - 2022 : Centre de Recherche en Astrophysique du Qu\'ebec \\
%{\bf NASA K2 Grant} 2015 - 2019 : Observing Uranus with K2 \\
%{\bf NASA ADAP} 2015 - 2019 : Microlensing in the Kepler FOV \\
%{\bf NASA Kepler PSP} 2014 - 2019 : Uniform Modeling of Kepler's Planetary Candidates 
%{\bf NASA Kepler PSP} 2011 - 2014 : Uniform Modeling of Kepler's Planetary Candidates 

\vspace*{-.2in}

\section{\sc Space Astronomy Leadership}
{\bf Canadian POET Mission}: Principle Investigator \\
{\bf NEOSSat Mission}: Science Team, Software Development \\
{\bf Pandora Mission}: Core Science Team \\
%{\bf BRITE Mission}: Science Team \\
{\bf Kepler Mission}: Science Team, Science Office

\vspace*{-.1in}

\section{\sc Current HQP}
{\bf Undergraduate} : 2 (S. Liu, A. Roy) ; {\bf Graduate} : 3 (M. Lightheat, D. Lizotte, M. Matesic) ; % {\bf PDF} : 1 (J. Sikora);

\vspace*{-.1in}

\section{\sc Awards}
{\bf Emerging Scholar Award: June 2019}\\
\vspace*{-.15in}
%\begin{list2}
%\item To recognize outstanding researchers and creators whose work has made a significant impact on their disciplines or fields, June 2019
%\end{list2}
\vspace*{-.15in}
{\bf Nasa Exceptional Scientific Achievement Medals: August 2011, August 2016}\\
%\vspace*{-.15in}
%\begin{list2}
%\item Development of a Model to Integrate the Diverse Types of Kepler Measurements to Derive the Characteristics of the Planet and the Host Star, 
%\item Validation of 800+ Transiting Extrasolar Planets, August 2016
%\end{list2}

\vspace*{.1in}

\section{\sc Selected Refereed Papers\\
\ \\
412 papers \\
H$_{index}$ = 92\\
cites = 32 020 \\
}

%\vspace*{-.15in}


Lissauer, J.J., \& {\bf Rowe, J.F.} et al. 2024, PSJ, Updated Catalog of Kepler Planet Candidates: 

\vspace*{-.15in}
{\bf Rowe, J.F.} et al.. 2022, SPIE, 12180, id. 121800A, The POET mission: a Canadian space telescope

\vspace*{-.15in}
Gilbert, E.A., et al. 2020, AJ 160, 116, The First Habitable-zone Earth-sized Planet from TESS.

%\vspace*{-.15in}
%Burke, C. et al. 2019 AJ, 157, 143, Re-evaluating Small Long-period Confirmed Planets from Kepler
%\vspace*{-.15in}
%{\bf Rowe, J.F.} et al.. 2017, AJ, 153, 149, Time-series Analysis of Broadband Photometry of Neptune

\vspace*{-.15in}
{\bf Rowe, J.F.} et al.. 2015, ApJS, 217, 16 Planetary Candidates Observed by Kepler, V: Planet Sample% from Q1-Q12 (36 Months)

\vspace*{-.15in}
{\bf Rowe, J.F.} et al.. 2014, ApJ, 784, 45 Validation of Kepler's Multiple Planet Candidates. III. Light %Curve Analysis and Announcement of Hundreds of New Multi-planet Systems

\end{resume}
\end{document}


%\vspace*{-.15in}
Kostov et al. 2021, AJ Accepted, TIC 172900988: A Transiting Circumbinary Planet Detected in One Sector of TESS Data

%\vspace*{-.15in}
Jontof-Hutter et al. 2021, AJ, 161, 246, Following Up the Kepler Field: Masses of Targets for Transit Timing and Atmospheric Characterization

%\vspace*{-.15in}
Gilbert, E.A., et al. 2020, AJ 160, 116, The First Habitable-zone Earth-sized Planet from TESS. I. Validation of the TOI-700 System

%\vspace*{-.15in}
Sikora, J., et al. 2020, MNRAS, 496, 295, Spectropolarimetric follow-up of 8 rapidly rotating, X-ray bright FK Comae candidates

%\vspace*{-.15in}
Kostov et al. 2020, AJ, 159, 253, TOI-1338: TESS' First Transiting Circumbinary Planet

%\vspace*{-.15in}
Sulis, S, A\&A, 631, 129, Multi-season optical modulation phased with the orbit of the super-Earth 55 Cancri e

%\vspace*{-.15in}
Kostov et al., 2019, AJ, 158, 32 The L 98-59 System: Three Transiting, Terrestrial-size Planets Orbiting a Nearby M Dwarf

%\vspace*{-.15in}
Burke, C. et al. 2019 AJ, 157, 143, Re-evaluating Small Long-period Confirmed Planets from Kepler

%\vspace*{-.15in}
Li, J.  et al., 2019, PASP, 131, 4506, Kepler Data Validation II-Transit Model Fitting and Multiple-planet Search

%\vspace*{-.15in}
Weiss, L. et al., 2018, AJ, 156, 254, The California-Kepler Survey. VI. Kepler Multis and Singles Have Similar Planet and Stellar Properties Indicating a Common Origin

%\vspace*{-.15in}
Mullally, F., et al., 2018, AJ, 155, 210,   Kepler’s Earth-like Planets Should Not Be Confirmed without Independent Detection: The Case of Kepler-452b

%\vspace*{-.15in}
Thompson et al. 2018, ApJS, 235, 38, Planetary Candidates Observed by Kepler. VIII. A Fully Automated Catalog with Measured Completeness and Reliability Based on Data Release 25 

%\vspace*{-.15in}
{\bf Rowe, J.F.} et al.. 2017, AJ, 153, 149, Time-series Analysis of Broadband Photometry of Neptune

%\vspace*{-.15in}
{\bf Rowe, J.F.} et al.. 2015, ApJS, 217, 16 Planetary Candidates Observed by Kepler, V: Planet Sample from Q1-Q12 (36 Months)

%\vspace*{-.15in}
{\bf Rowe, J.F.} et al.. 2014, ApJ, 784, 45 Validation of Kepler's Multiple Planet Candidates. III. Light Curve Analysis and Announcement of Hundreds of New Multi-planet Systems

%\vspace*{-.15in}
\textbf{Rowe J.F.} et al. 2010, ApJ, 713, 150, Kepler Observations of Hot Compact Objects

%\vspace*{-.15in}
\textbf{Rowe J.F.} et al. 2008, ApJ, 689, 1345 The Very Low Albedo of an Extrasolar Planet%: MOST Space-based Photometry of HD 209458




